\scanpage{120}
\mspage{35}
% archive sheet number: 119

% L01-03
jointe à 4.9, \inserted{qui implique que $E_2^{pq}=0$ si $q\ne0$,} donne l'isomorphisme voulu d'anneaux gradués :
\[
H^*(U^\infty/H,\Lambda)\simeq H^*(H,\Lambda).
\]

% L04-09
D'autre part, cet isom transforme, comme il est bien connu, les opérations de $G\deleted{(\unsure{\infty})}/H$ sur le premier mbr en les opérations déduites des opérations de $G\deleted{(\infty)}$ sur $H$ par autom. intérieurs. Comme $G\deleted{(\infty)}$ est commutatif, cette action est triviale, ce qui achève la démonstration de 4.10.

% L10-14
% relocation arrow from crossed Corollaire 4.12 toward the position before Remarque 4.12
\margin{\figure{\begin{tikzpicture}[scale=.45]\draw[->] (0,-2) to[out=150,in=220] (0,2);\end{tikzpicture}}}
\uline{Remarque} 4.12. Si on suppose seulement que les $X(m)$ satisfont au th. de pureté coh. en dim $\leq N$, alors les arguments qui précèdent prouvent encore 4.9 à 4.11 en dim $\leq N$, \deleted{et par suite aussi le corollaire 4.13 qui suit.}

% L15-18
\uline{Corollaire} 4.11. \deleted{4.12.} \inserted{Soit $\Gamma'$ l'image de $G\deleted{(\infty)}$ dans $\Gamma$, et} Supposons $\Gamma$ commutatif (ou plus gén. \deleted{l'image de $G(\infty)$ dans $\Gamma$}\inserted{$\Gamma'$ invariant dans $\Gamma$})\deleted{\unsure{invariant}}. Alors, l'action de $\Gamma$ sur $H^*(U',\Lambda)$ induit sur $\Gamma'$ l'action triviale (plus précisément $\Gamma'$ est le s-groupe des $g\in\Gamma$ qui opèrent trivialement sur $H^*(U',\Lambda)$).

% L19-24
Cet énoncé suffira \deleted{\unsure{pour le}} pour prouver \deleted{4.13. Nous \unsure{prouverons} maintenant} \deleted{plus bas} plus bas le th. de monodromie sous la forme 3.2. Avant de le faire, explicitons le calcul de $H^*(U',\Lambda)$ résultant de (4.10.1), ce qui nous sera utile dans l'exposé suivant.
